\documentclass[11pt,a4paper]{article}
% Common preamble for the hint sheets (\input from each file)
\usepackage[utf8]{inputenc}
\usepackage[T1]{fontenc}
\usepackage{lmodern}
\usepackage{amsmath,amssymb,amsthm,mathtools}
\usepackage[a4paper,margin=2.3cm]{geometry}
\usepackage{enumitem}
\usepackage{xcolor}
\usepackage[hidelinks]{hyperref}
\renewcommand{\leq}{\leqslant}
\renewcommand{\le}{\leqslant}
\renewcommand{\geq}{\geqslant}
\renewcommand{\ge}{\geqslant}
\setlength{\parindent}{0pt}
\setlength{\parskip}{0.35em}
% \hintstitle{Sheet title}{Class work / Homework}
\newcommand{\hintstitle}[2]{\begin{center}{\LARGE\bfseries Hints: #1\par}\vspace{0.5em}{\large #2\par}\end{center}
\noindent\textit{The hints to each problem are cumulative. Read one, try again, and open the next only if you are still stuck.}\par\medskip\hrule\medskip}
\newcommand{\problem}[1]{\par\bigskip{\large\bfseries Problem #1}\par\nopagebreak\smallskip}
\newcommand{\subpart}[1]{\par\smallskip\textbf{(#1)}\par\nopagebreak}
\newcommand{\hint}[1]{\par\smallskip\noindent\textcolor{black!70}{\textbf{Hint #1.}}\enspace}
\newcommand{\note}[1]{\par\smallskip\noindent\textit{Note.} #1\par}
\providecommand{\lt}{<}
\providecommand{\gt}{>}
\providecommand{\tightlist}{\setlength{\itemsep}{0pt}\setlength{\parskip}{0pt}}
\emergencystretch=3em

\begin{document}
\hintstitle{Abstract algebra -- 1. Finite groups: counting}{Homework}

\problem{1}
\hint{1} In the commutative case, raise $ab$ to the power $rs$: $(ab)^{rs}=a^{rs}b^{rs}=e$. Together with $(ab)^t=e$ and $\gcd(rs,t)=1$ this gives $ab=e$.
\hint{2} So $b=a^{-1}$, hence $a^r=e$ and $a^s=b^{-s}=e$. Since $\gcd(r,s)=1$, $a=e$.
\hint{3} In general the answer is no. Look for $a$ of order $2$ and $b$ of order $3$ with $ab$ of order $5$, e.g.\ in $A_5$: $a=(12)(34)$, $b=(135)$ (check that $ab$ is a $5$-cycle). Then $r=2$, $s=3$, $t=5$ work.

\problem{2}
\hint{1} Every $x\ne e$ commutes with its powers. If $\operatorname{ord}(x)\ge4$, then $x,x^2,x^{-1}$ give two different commuting pairs. So all non-identity elements have order $2$ or $3$. If there were two different subgroups of order $3$, they would give two pairs.
\hint{2} If no element has order $3$, then $x^2=e$ for all $x$, $G$ is abelian, and the number of commuting pairs is never exactly one. So there is a unique subgroup $C=\{e,c,c^2\}$ of order $3$, and $\{c,c^2\}$ is the only commuting pair. Hence $C(c)=C$, and $C$ is normal.
\hint{3} Conjugation gives a homomorphism $G\to\operatorname{Aut}(C)\cong\mathbb Z/2$ with kernel $C(c)=C$. So $|G|\in\{3,6\}$. A group of order $6$ in which $c$ commutes with no involution is non-abelian, hence $S_3$.

\problem{3}
\subpart{a}
\hint{1} $|H\cap xH|=|H|+|xH|-|H\cup xH|\ge 2p-n$.
\subpart{b}
\hint{1} Take $x\in H$. If $y=xh\in H\cap xH$ with $h\in H$, then $x,h,xh$ are all involutions (or $e$). Show that this forces $xh=hx$.
\hint{2} So the centraliser $C(x)$ contains the $\ge 2p-n>n/2$ elements $h\in H$ with $xh\in H$. A subgroup with more than $n/2$ elements is all of $G$, so every $x\in H$ is central.
\hint{3} The subgroup generated by $H$ has more than $n/2$ elements, so it is $G$. A group generated by central elements is abelian.
\subpart{c}
\hint{1} If $G$ were abelian, $H$ would be a subgroup. With $|H|>n/2$ this forces $H=G$, i.e.\ $p=n>3n/4$, a contradiction.

\problem{4}
\hint{1} Both statements are true. For (b), try groups with very few subgroups.
\hint{2} (b) $\mathbb Z/p$ with $p$ prime has exactly $2$ subgroups, so $n/s=p/2$. (a) Take $G=(\mathbb Z/2)^k$, an $\mathbb F_2$-vector space: its subgroups are its subspaces. Count the subspaces of dimension $\lfloor k/2\rfloor$, or just bound their number from below by about $2^{k^2/4}$, which is much larger than $2^k$.

\problem{5}
\note{The ``only if'' part fails for $n=1$, since then all $H^x=\{e\}$. The hints assume $n\ge2$; the formula for the union holds in all cases.}
\hint{1} If $x^{-1}y\in H$, then $H^y=H^{x(x^{-1}y)}=H^x$. Conversely, $H^x=H^y$ gives $H^{x^{-1}y}=H$, and by hypothesis this is impossible if $x^{-1}y\notin H$ (it would force $H=\{e\}$).
\hint{2} So the distinct conjugates correspond to left cosets of $H$: there are $m/n$ of them. Two distinct ones meet only in $e$, since $H^x\cap H^y=(H\cap H^{x^{-1}y})^x$.
\hint{3} Hence $\bigl|\bigcup_x H^x\bigr|=\frac mn(n-1)+1=m-\frac mn+1$.

\problem{6}
\hint{1} First produce a neutral element. Choose $e$ with $c*e=c$, and show that $w*e=w$ for every $w$: write $w=c*u$, using the hypothesis with $x=c$.
\hint{2} Then pick $c'$ with $c*c'=e$ and multiply $a*c=b*c$ by $c'$, using associativity and commutativity.

\problem{7}
\hint{1} If $G/Z(G)=\langle gZ\rangle$, write $x=g^iz_1$, $y=g^jz_2$ with $z_1,z_2$ central and multiply out.
\hint{2} For $|G|=p^2$: the class equation gives $Z(G)\ne\{e\}$, so $G/Z(G)$ has order $1$ or $p$ and is cyclic.
\hint{3} For order $p^3$, take the upper unitriangular $3\times3$ matrices over $\mathbb F_p$ (the Heisenberg group). For $p=2$ this is the dihedral group of order $8$.

\problem{8}
\hint{1} Let $K=\ker\pi$ and take $H\le G$ \emph{minimal} with $\pi(H)=Q$. Note $|H|=|H\cap K|\cdot|Q|$. Suppose a prime $r$ divides $|H|$ but not $|Q|$, and let $P$ be a Sylow $r$-subgroup of $H\cap K$. Then $P$ is even a Sylow $r$-subgroup of $H$.
\hint{2} Frattini argument: $H\cap K\trianglelefteq H$, and all Sylow $r$-subgroups of $H\cap K$ are conjugate inside $H\cap K$, so $H=(H\cap K)\,N_H(P)$. Hence $\pi(N_H(P))=Q$, and minimality gives $N_H(P)=H$, i.e.\ $P\trianglelefteq H$.
\hint{3} A normal Sylow subgroup has a complement $C$, with $H=PC$ and $|C|=|H|/|P|$ (Schur--Zassenhaus; this is the one tool beyond Sylow). Then $\pi(C)=\pi(PC)=Q$ because $P\subseteq K$, and $C<H$, which contradicts minimality.

\problem{9}
\subpart{a}
\hint{1} Let $x$ have order $d$. Show that $C(x)=\langle x\rangle$. If not, some prime $q$ divides $|C(x)|/d$, and $q\nmid d$ because $n$ is squarefree. Use Cauchy in $C(x)$ to get an element $y$ of order $q$ commuting with $x$, and look at the order of $xy$.
\hint{2} So the conjugacy class of every element of order $d$ has exactly $n/d$ elements, and the set of elements of order $d$ is a union of such classes.
\subpart{b}
\hint{1} By (a), $N\ge n/d$. The generators of $\langle x\rangle$ give $N\ge\varphi(d)$. Multiply, and use that $n/d$ and $d$ are coprime: $N^2\ge \frac nd\varphi(d)\ge\varphi(n/d)\varphi(d)=\varphi(n)$.

\problem{10}
\hint{1} Only inner automorphisms are needed. For each $g$ there is $m=m_g$ with $gxg^{-1}=x^{m}$ for all $x$. Apply this to $x=g$: you get $\operatorname{ord}(g)\mid m_g-1$.
\hint{2} Then $[g,y]=gyg^{-1}y^{-1}=y^{m_g-1}$ lies in $\langle y^{\operatorname{ord}(g)}\rangle$. By symmetry $[y,g]=[g,y]^{-1}$ lies in $\langle g^{\operatorname{ord}(y)}\rangle$.
\hint{3} Let $a=\operatorname{ord}(g)$, $b=\operatorname{ord}(y)$, $\delta=\gcd(a,b)$. The order of $[g,y]$ divides both $b/\delta$ and $a/\delta$, which are coprime. Hence $[g,y]=e$.

\end{document}
